
\medskip
\phantomsection\label{sec:sga3-VIA-2.5}%
\noindent\SGAInlineRefTarget{sga3:vi:paragraph:2-5}\textbf{2.5.}
This paragraph has been added. The results that follow appear in
Expos\'e~VI\textsubscript{B}, but could (or should) have appeared in
VI\textsubscript{A}; it is useful to have them available already, in
order to make \hyperref[thm:sga3-VIA-3.2]{Theorem~3.2} below more precise.

\medskip
\phantomsection\label{lem:sga3-VIA-2.5.1}%
\noindent\textbf{Lemma 2.5.1.}
\textit{Let \((A,\mathfrak m)\) be an Artinian local ring, and let
\(k=A/\mathfrak m\) be its residue field.}
\begin{enumerate}
\item[\SGAInlineRefTarget{lem:sga3-VIA-2.5.1:item:i}(i)] \textit{If \(X\) is an \(A\)-scheme such that
\(X\otimes_Ak\) is locally of finite type (respectively, of finite type)
over \(k\), then \(X\) has the same property over \(A\).}
\item[\SGAInlineRefTarget{lem:sga3-VIA-2.5.1:item:ii}(ii)] \textit{Let \(u:X\to Y\) be a morphism of \(A\)-schemes. If
\(u\otimes_Ak\) is an immersion (respectively, a closed immersion), then
so is \(u\).}
\end{enumerate}

\noindent\textit{Proof.}
\SGARefLink{lem:sga3-VIA-2.5.1:item:i}{(i)} Suppose that \(X\otimes_Ak\) is locally of finite type over \(k\).
Let \(U=\operatorname{Spec}B\) be an affine open subset of \(X\). By
hypothesis, there are elements \(x_1,\ldots,x_n\) of \(B\) whose images
generate \(B/\mathfrak mB\) as a \(k\)-algebra, and the ``nilpotent
Nakayama lemma'' then shows that the \(x_i\) generate \(B\) as an
\(A\)-algebra. This proves that \(X\) is locally of finite type over
\(A\). If, in addition, \(X\otimes_Ak\) is quasi-compact, then so is
\(X\), which has the same underlying topological space; hence \(X\) is
of finite type over \(A\). This proves \SGARefLink{lem:sga3-VIA-2.5.1:item:i}{(i)}.

We prove \SGARefLink{lem:sga3-VIA-2.5.1:item:ii}{(ii)}. Suppose that \(u\otimes_Ak\) is an immersion
(respectively, a closed immersion). Then \(u\) is a homeomorphism of
\(X\) onto a locally closed (respectively, closed) subset of \(Y\), and,
for every \(x\in X\), the ring homomorphism
\[
  \varphi_x:\mathcal O_{Y,u(x)}\longrightarrow\mathcal O_{X,x}
\]
is such that \(\varphi_x\otimes_Ak\) is surjective. The nilpotent
Nakayama lemma therefore shows that \(\varphi_x\) is surjective, so
\(u\) is an immersion (respectively, a closed immersion).

\medskip
\phantomsection\label{prop:sga3-VIA-2.5.2}%
\noindent\textbf{Proposition 2.5.2.}
\textit{Let \(A\) be an Artinian local ring with residue field \(k\), and
let \(u:G\to H\) be a quasi-compact morphism between \(A\)-groups locally
of finite type.}
\begin{enumerate}
\item[\SGAInlineRefTarget{prop:sga3-VIA-2.5.2:item:a}(a)] \textit{The set \(u(G)\) is closed in \(H\), and its connected
components are irreducible and all have the same dimension.}
\item[\SGAInlineRefTarget{prop:sga3-VIA-2.5.2:item:b}(b)] \textit{One has}
\[
  \dim G=\dim u(G)+\dim\operatorname{Ker}(u).
\]
\item[(c)] \textit{If \(u\) is a monomorphism, then it is a closed
immersion.}
\end{enumerate}

\noindent\textit{Proof.}
By \hyperref[lem:sga3-VIA-2.5.1]{the preceding lemma}, it suffices to
prove the proposition when \(A=k\). Moreover, the properties under
consideration are stable under fpqc descent, and dimension is invariant
under extension of the base field, so we may assume that \(k\) is
algebraically closed.

We prove \SGARefLink{prop:sga3-VIA-2.5.2:item:a}{(a)}. Let \(C\) be the reduced subscheme of \(H\) whose
underlying topological space is \(u(G)\). Since \(u(G)\) is stable under
the inversion morphism of \(H\), so is \(C\). Furthermore,
\(u:G\to C\) is quasi-compact and dominant. Thus, by
EGA~IV\textsubscript{2}, 2.3.7, the same is true of
\[
  u\times_k\operatorname{id}_G
  \qquad\text{and}\qquad
  \operatorname{id}_H\times_ku,
\]
and hence of their composite
\[
  u\times_ku:G\times_kG\longrightarrow C\times_kC.
\]
Consequently, multiplication on \(H\) sends \(C\times_kC\) into \(C\),
and therefore \(C\) is a group subscheme of \(H\).

Replacing \(H\) by \(C\), we are thus reduced to the case in which \(u\)
is dominant. Then \(G(k)\) is dense in \(H\), and hence meets every
connected component of \(H\); it therefore acts transitively on the set
of those connected components. It is consequently enough to show that
\(u(G)\) contains \(H^0\). Replacing \(G\) by \(u^{-1}(H^0)\), we may
therefore suppose that \(H=H^0\). In that case,
\hyperref[prop:sga3-VIA-2.4]{2.4} shows that \(H\) is irreducible and of
finite type over \(k\), and hence Noetherian. Moreover, \(u\) is locally
of finite type (cf.\ EGA~I, 6.6.6) and quasi-compact, and is therefore of
finite type. By Chevalley's constructibility theorem
(cf.\ EGA~IV\textsubscript{1}, 1.8.5), \(u(G)\) is consequently a
constructible (and dense) subset of
\[
  H=\overline{u(G)},
\]
and thus contains a dense open subset \(U\) of \(H\)
(cf.\ EGA~0\textsubscript{III}, 9.2.2). By
\hyperref[sga3:VIA:0.5]{0.5}, we then have
\[
  H=U\cdot U\subset u(G),
\]
whence \(u(G)=H\). Together with
\hyperref[cor:sga3-VIA-2.4.1]{2.4.1}, this proves \SGARefLink{prop:sga3-VIA-2.5.2:item:a}{(a)}.

We prove \SGARefLink{prop:sga3-VIA-2.5.2:item:b}{(b)}. Recall first that the functor
\(\operatorname{Ker}(u)\) (cf.\ I, \SGARefLink{sga3:i:definition:2-3-6-1}{2.3.6.1}) is represented by
\(u^{-1}(e)\), where \(e\) denotes the identity element of \(H\). Since
\(u\) is of finite type, \(\operatorname{Ker}(u)\) is of finite type
over \(k\). On the other hand, replacing \(H\) by the reduced closed
subscheme \(u(G)\), we may assume that \(u\) is surjective. Let \(u^0\)
denote the restriction of \(u\) to \(G^0\). Since \(G\) and
\(\operatorname{Ker}(u)\) are equidimensional, and since
\[
  \operatorname{Ker}(u)^0\subset\operatorname{Ker}(u^0),
\]
we are reduced to the case in which \(G\), and hence also \(H\), is
irreducible.

By EGA~IV\textsubscript{3}, 9.2.6.2 and 10.6.1(ii), the set of
\(y\in H\) such that
\[
  \dim u^{-1}(y)=\dim G-\dim H
\]
contains a nonempty open subset \(V\). Since \(u\) is surjective,
\(U=u^{-1}(V)\) is then a nonempty open subset of \(G\), and therefore
contains a closed point \(x\) of \(G\), because \(G\) is a Jacobson
scheme (cf.\ EGA~IV\textsubscript{3}, 10.4.8). Right translation \(r_x\)
is an isomorphism from \(\operatorname{Ker}(u)\) onto
\(u^{-1}(u(x))\), whence
\[
  \dim\operatorname{Ker}(u)
    =\dim u^{-1}(u(x))
    =\dim G-\dim H.
\]

We prove (c), following \textup{[DG70]}, I, \S~3.4. (Another proof is
given in Expos\'e~VI\textsubscript{B}, 1.4.2.) Suppose that \(u\) is a
monomorphism. If \(C\) is a connected component of \(G\), there is a
closed point \(x\in G\) such that \(C=r_x(G^0)\). If \(u_C\)
(respectively, \(u^0\)) denotes the restriction of \(u\) to \(C\)
(respectively, to \(G^0\)), then
\[
  u_C=r_{u(x)}\circ u^0\circ r_x^{-1}.
\]
It therefore suffices to show that \(u^0\) is a closed immersion. We may
thus suppose that \(G=G^0\), so that \(G\) is irreducible and of finite
type over \(k\).

Let \(\xi\) be the generic point of \(G\). Then
\(\mathcal O_{G,\xi}\) is an Artinian local ring; denote its maximal
ideal by \(\mathfrak m\). Let \(h=u(\xi)\), let \(\mathfrak n\) be the
maximal ideal of \(\mathcal O_{H,h}\), and set
\[
  A=\mathcal O_{G,\xi}/\mathfrak n\mathcal O_{G,\xi}.
\]
Since \(u\) is a monomorphism, so is the morphism
\[
  u_h:\operatorname{Spec}(A)\longrightarrow
      \operatorname{Spec}(\kappa(h))
\]
obtained by base change. The multiplication morphism
\[
  A\otimes_{\kappa(h)}A\longrightarrow A
\]
is therefore an isomorphism (cf.\ EGA~I, 5.3.8), whence
\(A=\kappa(h)\). Nakayama's lemma then shows that
\[
  \mathcal O_{H,h}\longrightarrow\mathcal O_{G,\xi}
\]
is surjective, since
\(\mathfrak n\mathcal O_{G,\xi}\) is contained in \(\mathfrak m\) and
is therefore nilpotent.

Let \(V\) be an affine open subset of \(H\) containing \(h\), let
\(U\ne\varnothing\) be an affine open subset of \(G\) contained in
\(u^{-1}(V)\), and let
\[
  \varphi:\mathcal O_H(V)\longrightarrow\mathcal O_G(U)
\]
be the homomorphism of \(k\)-algebras induced by \(u\). Let
\(\mathfrak p\) be the prime ideal of \(\mathcal O_G(U)\) corresponding
to \(\xi\),\begingroup
\renewcommand{\thefootnote}{\fnsymbol{footnote}}%
\footnote[1]{Translator's note: the French source prints \(\eta\) here;
the generic point introduced immediately above is \(\xi\).}%
\endgroup
and set
\(\mathfrak q=\varphi^{-1}(\mathfrak p)\). Since
\(G\) is of finite type over \(k\), the \(k\)-algebra
\(\mathcal O_G(U)\) is generated by finitely many elements
\(a_1,\ldots,a_n\). By the preceding argument, there are elements
\(b_1,\ldots,b_n,s\in\mathcal O_H(V)\), with
\(s\notin\mathfrak q\), such that in
\(\mathcal O_G(U)_{\mathfrak p}\) one has
\[
  a_i/1=\varphi(b_i)/\varphi(s).
\]
There are therefore elements
\[
  t_1,\ldots,t_n\in\mathcal O_G(U)-\mathfrak p
\]
such that
\[
  t_i\bigl(a_i\varphi(s)-\varphi(b_i)\bigr)=0.
\]
Set
\[
  t=t_1\cdots t_n\varphi(s)\in\mathcal O_G(U)-\mathfrak p.
\]
The equalities \(a_i/1=\varphi(b_i)/\varphi(s)\) then already hold in
\(\mathcal O_G(U)_t\). Since
\[
  t\in\mathcal O_G(U)=k[a_1,\ldots,a_n],
\]
there is an element \(b\in\mathcal O_H(V)\) such that, for some
\(r\in\mathbb N\),
\[
  t/1=\varphi(b)/\varphi(s)^r.
\]
Consequently, \(\varphi\) induces a surjection
\[
  \mathcal O_H(V)_{sb}\longrightarrow\mathcal O_G(U)_t,
\]
and hence \(u\) is a local immersion at \(\xi\).

The open subset \(W\) of \(G\) consisting of the points at which \(u\)
is a local immersion is therefore nonempty. Since \(G\) is a Jacobson
scheme, \(W\) contains a closed point \(y\). To prove that \(W=G\), it
is enough to show that every closed point of \(G\) belongs to \(W\). But
every closed point \(x\) is the image of \(y\) under the translation
\(r_x\circ r_y^{-1}\), and hence belongs to \(W\). Thus \(W=G\), which
proves that \(u\) is a local immersion.

Since \(G\) is irreducible, it follows that \(u\) is an immersion.
Indeed, for every \(x\in G\), let \(U_x\) and \(V_x\) be open subsets of
\(G\) and \(H\), respectively, such that \(x\in U_x\) and \(u\)
induces a closed immersion of \(U_x\) into \(V_x\). Since \(U_x\) is
dense in \(G\), the subset \(u(U_x)\) is dense in \(u(G)\cap V_x\).
But \(u(U_x)\) is closed in \(V_x\), so
\footnote{Translator's note: the French re-edition prints
\(u(U_x)=V_x\). The preceding density statement and the next inverse-image
identity force the intersection shown here.}
\[
  u(U_x)=u(G)\cap V_x.
\]
Moreover, since \(u\) is injective, \(U_x=u^{-1}(V_x)\). Thus \(u\)
induces a closed immersion of \(G\) into the open subscheme of \(H\)
covered by the \(V_x\). Hence \(u:G\to H\) is an immersion. We have
already seen that \(u(G)\) is closed in \(H\), so \(u\) is a closed
immersion.
